The channel-by-channel ablation showed that detection performance was concentrated specifically at the frontopolar electrodes rather than across the frontal region as a whole (Table~\ref{tab:exp1_channel_ablation}). In Raja, Fp2 (EGI E9) achieved $F_1=0.868$ and Fp1 (EGI E22) achieved $F_1=0.863$, whereas performance decreased at AF4 and AF3 to $F_1=0.803$ and $F_1=0.738$, respectively, and declined further at F3 and F4 to $F_1=0.456$ and $F_1=0.452$. The same localisation was observed in Cao2018, where Fp1 and Fp2 achieved $F_1=0.797$ and $F_1=0.784$, compared with $F_1=0.595$ at F3 and $F_1=0.580$ at F4. Performance collapsed outside frontal sites, with the worst electrodes reaching $F_1=0.017$ in Raja and $F_1=0.028$ in Cao2018. Despite the different montages, both corpora therefore localised the optimum to the corresponding Fp1/Fp2 frontopolar sites, with the Raja EGI equivalents listed in Table~\ref{tab:egi_map}. The best single frontopolar electrode approached the whole-cap oracle values of $F_1=0.8894$ and $F_1=0.8175$, indicating that a dense montage was not required.
